\section{Resultados}\label{chap:resultados-final}

\subsection{Artefato desenvolvido}

A plataforma integra seleção e reutilização de geometrias, processamento Sentinel-2, exploração espacial, séries, contexto meteorológico e resultados portáveis. O processamento remoto evita manter o acervo orbital no servidor da aplicação; o pré-carregamento raster permite consulta numérica no navegador. Armazenamento local e isolamento por sessão atendem a contextos distintos, complementados por transferência GeoJSON.

\subsection{Avaliação computacional}

A versão avaliada corresponde à base \texttt{b7c325c} acrescida das alterações locais documentadas. Em Linux x86\_64, Python 3.14.4, pytest 9.0.3 e tracer C, o comando \texttt{python -m pytest -q -p no:cacheprovider}, com \texttt{COVERAGE\_CORE=ctrace}, aprovou \textbf{710 testes em 11,73 s}, sem falhas ou testes ignorados. A cobertura em \texttt{src/app} foi de \textbf{100\% das linhas e ramificações}, com 2.735 instruções e 856 ramificações, sem caminhos parcialmente cobertos. A linha de base de 686 testes permaneceu preservada; novos casos verificaram os contratos corrigidos.

Os testes abrangem regras locais, integração entre módulos, persistência temporária, widgets Streamlit e serialização. Serviços remotos são substituídos em parte da suíte. A medição não inclui o ponto de entrada da raiz nem HTML/JavaScript. Separadamente, três testes Node verificaram legenda negativa, lookup de células e NoData, e descompressão zlib/Base64 em \texttt{Float32Array}; os três foram aprovados. Esses testes não constituem uma matriz de compatibilidade de navegadores.

Cobertura integral indica execução dos caminhos medidos, não ausência de defeitos. A avaliação dos contratos encontrou suposições semânticas que a suíte anterior não distinguia, motivando regressões específicas sem reduzir a meta.

\subsection{Verificação das integrações}\label{sec:resultado-unidades}

Uma resposta pública NASA POWER, preservada para longitude $-52$ e latitude $-28$ entre 01/01/2025 e 02/01/2025, declara radiação em MJ/m²/dia e valores 21,64 e 20,99. A média é 21,315. A conversão anterior multiplicava essa média por 3,6, produzindo 76,734 sob o mesmo rótulo. A normalização pelos metadados passou a preservar MJ e converter somente kWh declarados. Testes verificaram a média corrigida, as duas linhas XLSX/CSV e a representação resumida do PDF.

A avaliação também harmonizou limites temporais, resumo meteorológico compartilhado e metadados dos pipelines. Um teste multiíndice confirmou uma única obtenção do contexto climático. Outro distinguiu aquisições do mesmo dia na marcação de anomalias. Alterações analíticas passaram a invalidar todos os produtos derivados, enquanto a troca visual entre índices calculados preserva a análise.

No mapa, o mesmo intervalo usado nos tiles é enviado à legenda. A escala fixa do PDF permanece uma escolha separada, comunicada no próprio produto. A correção do contrato não implica igualdade pixel a pixel entre reduções remotas, grade reamostrada e mapa estático.

\subsection{Demonstração funcional}

Uma sessão real em Chrome confirmou o fluxo entre busca por Passo Fundo, demarcação, salvamento, restauração, análise e download. O recorte demonstrativo de aproximadamente 4,36 ha foi consultado entre 01/09/2026 e 31/10/2026, sem atribuição de cultura ou propriedade. Duas imagens passaram pelos filtros, em 08/09 e 20/10. As médias espaciais das composições foram 0,7070 (NDVI), 0,2140 (NDWI) e 0,2380 (NDMI). Esses valores caracterizam o exemplo computacional, não uma amostra representativa de áreas agrícolas.

As séries contiveram dois pontos por índice, insuficientes para a janela de três observações anteriores do detector; a ausência de anomalias nessa demonstração não informa eficácia de detecção. O clima apresentou 61 registros, precipitação acumulada de 635,56 mm e temperatura média de 16,24 °C. A radiação permaneceu indisponível e foi mantida ausente nos arquivos, sem preenchimento por suposição.

O PDF baixado contém três páginas com mapas estáticos, o XLSX possui sete abas e o ZIP sete arquivos CSV com integridade estrutural conferida. O identificador de aquisição está presente nas séries tabulares. A consulta NDMI na posição $(-52{,}331,-28{,}228)$ exibiu 0,3228, tanto por movimento quanto por toque em viewport de 390 pixels. O resultado pertence à célula reamostrada da composição e não substitui a média espacial de 0,2380.

O endereço público respondeu HTTP 200 e apresentou a interface Streamlit em navegador. Essa verificação não estabelece equivalência entre versão publicada e alterações locais nem disponibilidade contínua. As capturas analíticas correspondem à execução local corrigida; manifestos e arquivos preservam o vínculo com os parâmetros.

\subsection{Avaliação com usuários}

O instrumento está disponível para participação voluntária, conforme a \cref{sec:pesquisa-usuarios}. Não foi disponibilizado um conjunto de respostas reais para esta análise; por isso, não são apresentados tamanho amostral, avaliações, gráficos ou depoimentos. O procedimento reproduzível foi preparado e testado com entradas sintéticas isoladas, que não integram os resultados da pesquisa.
